\section{Introduction}

Software systems built around AI agents are increasingly able to emit logs, traces, and runtime events at high volume. Those materials are useful for debugging, observability, and post hoc inspection. In FDO-based agent systems, however, a narrower and harder question often remains unresolved: for one concrete operation, which operation was performed on which digital object context, under which policy constraints, through which input and output references, with what evidence, and through what validation path? In other words, it is not enough to know that an execution happened. A third party may need a compact statement that can be packaged, checked, and archived as an accountability unit for a single operation.

That gap is the starting point of this paper. The problem addressed here is not general AI governance, broad compliance automation, or a full ecosystem mapping for all FDO variants. The problem is smaller and more concrete. Given a single operation in an FDO-based agent system, can we represent its accountability conditions in a minimal form that is still independently verifiable? The implementation studied in this paper addresses that question through a deliberately constrained engineering path: a minimal profile, a profile-aware validator, a small boundary-oriented example suite, a command-line entry point, a runnable demo, and a blinded review-package surface. The aim is not breadth. The aim is to turn operation accountability from a loose narrative into a testable software artifact.

The need for such a minimal artifact becomes clearer when common alternatives are considered. Ordinary logs can record timestamps, events, paths, and fragments of execution state, but they do not usually form a stable accountability object. They are often distributed across system components, shaped by runtime convenience rather than portability, and weak in explicit policy binding. Provenance-oriented representations can describe derivation links and object relationships more clearly, but provenance alone does not necessarily say under which governing policy the operation was performed, what the minimal evidence block is, or how an external validator should classify failures. Policy-only representations can state constraints and allowed actions, but they do not establish what actually happened for a particular operation, whether input and output references close locally, or whether the statement is independently verifiable. These alternatives are partial capability holders. The narrower claim of this paper is that operation accountability benefits from a minimal unit that binds these dimensions together rather than leaving them split across unrelated surfaces.

The paper therefore adopts a minimal-profile approach. The current implementation uses a versioned operation-accountability profile organized around five fixed parts: \texttt{operation}, \texttt{policy}, \texttt{provenance}, \texttt{evidence}, and \texttt{validation}. Around those core parts, the profile retains only the supporting context needed to make the statement checkable, such as \texttt{actor}, \texttt{subject}, \texttt{constraints}, \texttt{profile}, and \texttt{timestamp}. This design is intentionally restrictive. It does not attempt to absorb every metadata field that future governance workflows might request. It does not attempt to model large workflows or multi-agent coordination. Instead, it treats a single operation accountability statement as the primary engineering object and asks what minimum structure is needed for independent verification.

That emphasis on verification shapes the validator as well. In the current implementation, validation is not reduced to a schema pass. The validator performs staged checks over the profile and is exposed through a command-line surface. It first checks structural conformance against the expected shape. It then checks local reference closure. It continues with cross-field consistency among policy, provenance, and evidence links. Finally, it applies minimal integrity checks. The output is a validation report with machine-readable JSON, a human-readable summary, and explicit error codes. This is why the validator is described as profile-aware: it does not merely parse a JSON document; it enforces the identity, linkage, and consistency rules defined for this profile.

The current implementation provides grounded evidence for this claim in a form suitable for a methodology-plus-artifact paper. At the time of writing, the implementation contains a profile specification, a JSON Schema, two valid examples, and five invalid examples. The invalid examples are intentionally minimal: each breaks one primary rule so that failure categories remain readable and diagnosable. The two valid examples are also informative beyond simple pass cases. One represents a metadata enrichment context with a single input and a single output. The other represents a retention review context with two inputs and one decision output. Both pass under the same profile model and the same validator path. That does not establish broad cross-framework validation, but it does provide limited evidence for basic portability and a second-context validity check for the same minimal structure.

The artifact surface is broader than the profile files alone. The implementation includes tests that exercise the valid and invalid examples, the command-line path, and the validator entry points. It also includes a runnable single-path demo that proceeds from object load or creation, through profile precheck and operation execution, to evidence generation and final validation output. In the anonymous review version, public release and archive identifiers are withheld, but the frozen review package preserves the same method path. Those anchors do not prove superiority over alternative approaches, but they matter for a software engineering paper because they show that the work is packaged as a releasable, citable, and archivable artifact rather than a purely conceptual proposal.

\subsection{Contributions and Scope Boundary}

This paper makes four concrete contributions.
\begin{enumerate}
\item It defines a minimal verifiable profile for single-operation accountability in FDO-based agent systems, centered on \texttt{operation}, \texttt{policy}, \texttt{provenance}, \texttt{evidence}, and \texttt{validation}.
\item It presents a profile-aware validation approach that extends structural schema checking with reference closure, cross-field consistency, and explicit error reporting.
\item It provides a boundary-oriented example suite with two valid examples and five single-rule invalid examples, making both passing and failing conditions inspectable.
\item It packages these elements into a reproducible artifact package spanning specification, schema, examples, validator, command-line validation, tests, demo, and a frozen review-package surface.
\end{enumerate}

The scope of the paper is deliberately minimal. The paper does not claim a universal FDO standard. It does not claim full support for complex multi-agent orchestration. It does not claim a complete cryptographic trust infrastructure, non-repudiation, or organizational compliance automation. It also does not claim broad empirical validation across frameworks or implementations. The current evidence is narrower: a single profile model, a single reference implementation, two valid contexts, five invalid boundary cases, a validator path, a command-line surface, a runnable demo, and a frozen review-package surface. The argument of the paper should be read at that scale.

That restraint keeps the contribution aligned with a methodology-, validator-, and artifact-oriented software engineering paper. The implementation does not try to win by scope inflation. Instead, it closes a small loop end to end. The specification, JSON Schema, validator path, example suite, command-line surface, demo, and blinded artifact surface make that loop directly inspectable in the review version. The comparative argument is similarly restrained. In the paper, ordinary logs, provenance-only representations, and policy-only representations are not dismissed; they are treated as partial representations that do not by themselves form the same minimal accountability unit. The contribution is therefore not a sweeping survey claim, but a concrete engineering claim about what has been implemented as a reproducible, verifiable object.

The rest of the paper follows that structure. It first defines the problem and the design goals that motivate a deliberately minimal profile. It then presents the profile model, the validation model, and the reference validator. Next, it describes the artifact package and evaluates the current evidence through examples, tests, command-line behavior, demo execution, and blinded artifact packaging. Finally, it discusses limits, threats to validity, and the boundary between the current basic portability evidence and the broader claims that the present implementation does not yet support.
